%% guide-curling-phases — les quatre phases du lancer d'une pierre de curling.
%% Quatre vignettes identiques : la même pierre sur la glace, avec le seul élément
%% caractéristique de la phase.
%% RÈGLE : aucun vecteur poids ni réaction normale (c'est le bilan demandé par l'exercice),
%% aucune valeur numérique.
\documentclass[tikz,border=4pt]{standalone}
\usepackage{liaisons}
\begin{document}
\begin{tikzpicture}[x=1cm,y=1cm,
  glace/.style={line width=1.2pt, draw=solideB!40},
  sep/.style={line width=0.4pt, draw=black!25},
  vect/.style={-{Stealth[length=6pt]}, line width=1.2pt},
  leg/.style={font=\scriptsize, align=center, inner sep=1.5pt},
  phase/.style={font=\small\bfseries, anchor=north, inner sep=3pt}]

%% ---------- la pierre de curling, centrée en (#1,#2) ----------------------
\newcommand\pierre[2]{%
  \draw[glace] ({#1-0.85},{#2}) -- ({#1+0.85},{#2});
  \draw[line width=0.7pt, fill=black!30] ({#1-0.36},{#2+0.02}) rectangle ({#1+0.36},{#2+0.30});
  \draw[line width=0.7pt, fill=black!15] ({#1-0.30},{#2+0.30}) rectangle ({#1+0.30},{#2+0.44});
  \draw[line width=0.9pt] ({#1},{#2+0.44}) -- ({#1},{#2+0.58});
  \draw[line width=0.9pt] ({#1-0.12},{#2+0.58}) -- ({#1+0.12},{#2+0.58});}

\def\yg{0.40}   % niveau de la glace

%% ================= les quatre vignettes ===================================
\foreach \k/\lab/\leg in {0/A/{pierre posée}, 1/B/{poussée par le lanceur},
                          2/C/{glissement libre}, 3/D/{balayage} } {
  \pgfmathsetmacro\cx{2*\k+1}
  \pierre{\cx}{\yg}
  \node[phase] at (\cx,0.10) {\lab};
  \node[leg, anchor=south] at (\cx,1.92) {\leg}; }
\foreach \k in {1,2,3} { \draw[sep] (2*\k,-0.35) -- (2*\k,2.20); }

%% ================= A : la pierre est immobile =============================
\node[draw=black!55, line width=0.6pt, rounded corners=2pt, fill=white,
      font=\small, inner sep=2.5pt] at (1.00,1.35) {$\vec{v} = \vec{0}$};

%% ================= B : le lanceur pousse ==================================
\draw[line width=2.2pt, draw=black!45, line cap=round] (2.28,1.80) -- (2.62,1.12);
\draw[line width=0.7pt, draw=black!45, fill=black!20] (2.68,1.02) circle (0.10);
\draw[vect, draw=solideA] (3.02,0.98) -- (3.70,0.98);
\node[font=\small, text=solideA, anchor=south, inner sep=1.5pt] at (3.36,1.02) {$\vec{F}$};

%% ================= C : glissement libre ===================================
\draw[vect, draw=solideD] (4.60,1.30) -- (5.50,1.30);
\node[font=\small, text=solideD, anchor=south, inner sep=1.5pt] at (5.05,1.34) {$\vec{v}$};
\node[leg, text=black!65] at (5.00,0.98) {la vitesse diminue};

%% ================= D : balayage ===========================================
\foreach \bx in {7.44,7.72} {
  \draw[line width=0.9pt, draw=black!60] ({\bx+0.09},0.98) -- ({\bx+0.09},0.60);
  \hachures[draw=black!40]{({\bx-0.02},0.42) rectangle ({\bx+0.20},0.60)} }
\draw[vect, draw=solideD] (6.60,1.30) -- (7.40,1.30);
\node[font=\small, text=solideD, anchor=south, inner sep=1.5pt] at (7.00,1.34) {$\vec{v}$};
\node[leg, text=black!65] at (7.00,0.98) {vitesse constante};
\end{tikzpicture}
\end{document}
